\section{Approche réflexive}

\subsection{Si vous deviez ne retenir qu’une seule notion de ce travail, quelle serait-elle ?}
Les dysménorrhées doivent être explorées par le médecin. Trop souvent, les patientes traversent un désert diagnostique, ne se sentent pas soutenues ni écoutées. La première mission d’un médecin est de soigner un patient et non une pathologie. Apporter un maximum de confort aux patients est l'objectif.

\subsection{Que pensez-vous avoir bien réussi dans ce travail ?}
Je pense que ce cas clinique n’est pas un cas diagnostique idéal vu l’examen clinique pauvre et les diagnostics différentiels peu nombreux. D’autres cas auraient été plus évidents. Toutefois, cela m’a poussée à me renseigner davantage sur les différentes kystes ovariens possibles. Je pense, par contre, avoir gardé une structure claire tout au long du travail facilitant ainsi la compréhension. La recherche réalisée sur PubMed m’a grandement aidé pour appuyer la démarche diagnostique. Lire ces études m'a permis de mieux comprendre le contexte médical et les enjeux des différents examens cliniques possibles, leurs avantages et inconvénients. Être à jour au niveau de la littérature médicale est un avantage certain car elle permet un choix éclairé et optimal pour la patiente.

\subsection{Que pensez-vous devoir améliorer ?}
Je pense pouvoir améliorer la formulation de mes phrases, afin de les rendre plus médicalisées et précises. J’ai tenté au maximun d'utiliser le jargon médical, mais je suis convaincue qu’avec l’expérience, je parviendrai à m’exprimer de manière plus concise et rigoureuse. Lors de l’anamnèse, je ne suis pas toujours certaine de la pertinence des informations récoltées.

Le contexte socio-économique et les antécédents médicaux, notamment lorsqu’ils concernent un séjour à l’étranger, ont été peu développés dans ce travail en raison de données manquantes. La barrière de la langue représente également un obstacle à une communication optimale avec la patiente. Il aurait été pertinent d’obtenir davantage d’informations objectives sur les douleurs pelviennes rapportées depuis 7 ans.

\subsection{Quelle(s) question(s) souhaiteriez-vous approfondir dans un prochain travail ?}
L’aspect thérapeutique de la démarche clinique, particulièrement dans le cadre d’une pathologie chronique, est un point que je souhaiterais approfondir. Aujourd’hui encore, les liens entre lésions et symptomatologie ne sont pas toujours scientifiquement établis, ce qui explique les multiples options thérapeutiques proposées aux patientes.

J’aurais apprécié pouvoir suivre l’évolution de la patiente et comprendre son choix de s’orienter vers une solution chirurgicale. Ouvrir la discussion sur d’éventuelles alternatives aurait été enrichissant, tant pour elle que pour moi en tant que future médecin. Explorer ces dimensions aurait apporté une valeur ajoutée à ce travail.